\section{同时检测多个列表的成本}\label{sec:list}

两次互异性检测足以核查各条记录是否已抵达指定的目的地。这给定理~\ref{thm:many} 中的两种情形提供了共同的起点。对于无端口的域，各种可能的目的地迫使许多连线在适当谱空间中的总长度很大。对于有端口的有限域，如果能以很低的成本检测互异性，就能够压缩一个已知其下界的张量族。

\subsection{核查带标签的置换}

记 \(h=w-1\)、\(N=2^h\)，用 \(h\) 比特的字为记录编号。文献~\cite[定义~\ext{G}{def:tag}]{oaG} 中的带标签置换张量 \(\operatorname{Tag}(h,h)\) 的输入是 \(N\) 条记录 \((j_i,x_i)\) 以及含 \(N\) 个有效载荷 \(y_j\) 的有序列表，每个标签与每个有效载荷均占 \(h\) 比特。若 \(i\mapsto j_i\) 为置换且对每个 \(i\) 均有 \(x_i=y_{j_i}\)，张量便取值一。其余表项均为零。

第一次 \(D_w\) 检测用于核查标签。将每个 \(j_i\) 编码为 \((0,j_i,0^{h+1})\)，再在剩余的 \(N\) 个位置填入互异的字 \((1,j,0^{h+1})\)。这 \(2N\) 个字的长度均为 \(2w\)，而检测通过当且仅当标签构成一个置换 \(\pi\)。

第二次 \(D_w\) 检测用于核查有效载荷。其前 \(N\) 个字为 \((j_i,0,x_i,0)\)。在后半段中与 \(j\) 对应的位置，对两种选择求和：以权重 \(1\) 选取备用字 \((j,1,0^h,0)\)，以权重 \(-1\) 选取比较字 \((j,0,y_j,0)\)。如果标签构成 \(\pi\)，那么在地址 \(j\) 处，只有比较字有可能与来自 \(i=\pi^{-1}(j)\) 的字发生碰撞。因此，带符号的和可分解为
\[
     \prod_j \bigl(1-\mathbf 1_{\{x_{\pi^{-1}(j)}\ne y_j\}}\bigr)
       =\prod_j\mathbf 1_{\{x_{\pi^{-1}(j)}=y_j\}} .
\]
这一恒等式在任意特征下都成立。每个地址各有一个参与求和的选择比特，控制该处的字及其权重。复制标签、选取各字及固定常数共需 \(O(Nh)\) 的三价网络体积。因此，如果 \(D_w^{\otimes t}\) 有体积为 \(V\) 的实现，使用该实现两次，就能为 \(\operatorname{Tag}(h,h)^{\otimes t}\) 构造体积至多为 \(2V+O(tNh)\) 的网络，即使各列表在该实现中共用部分计算也成立。

\subsection{无端口的域}

设一类 Hadamard 空间容许一族阶数无界、且度数有共同上界的正则图，并满足以下性质：无论如何把图的顶点放置在该类中的任何一个空间中，独立均匀选取的两个顶点之间的平均平方距离总不超过边长平方的均值的某个固定常数倍。此时称该类空间\emph{有端口}。当一个域的 Hadamard 谱空间具有此性质时，就说该域有端口~\cite[定义~\ext{G}{def:ports}]{oaG}。

\begin{lemma}[均匀分布的目的地]\label{lem:uniform-demand}
设所用域无端口。考虑一个具有 \(r\) 个输入比特、\(r\) 个输出比特的固定三价网络拓扑。设对于某个分布的支撑中的每个置换，都能通过选取顶点取值表，实现将每个输入与其指定输出相连的 \(r\) 条独立连线。如果每个输入的目的地都在某个含 \(q\ge2\) 个输出的集合上均匀分布，则该拓扑的规模为 \(\Omega(r\sqrt{\log q})\)。目的地集合可以依输入而异。
\end{lemma}

\begin{proof}
给每个外部比特添上各自的相等关系叶节点。在用于度量距离的图上添加若干条边，将所得连通分支依次连成一条路径。这使最大度至多增加二，得到的边数记为 \(m\)，其中 \(m\) 至多是网络规模的一个常数倍。令 \(\alpha\) 表示均匀选取的一条有向边的分布，令 \(\mu\) 表示均匀选取的输入及其随机目的地的联合分布。对于度量为 \(d\) 的谱空间中的任意放置，都有
\[
                    r\,\mathbb E_\mu d\ \le\ m\,\mathbb E_\alpha d .
\]
事实上，对于每个置换，网络中的各张量通过局部收缩将内部比特连线的张量积变为指定的外部连线。由价格的单调性与可加性，各输入到其目的地的距离之和不超过内部边长之和。再对置换取平均，并在右侧计入所添加的边，即得所述不等式。

我们需要找到这样一种放置，使许多目的地同时很远。Hadamard 谱空间这一类在等距变换、缩放、取闭凸子空间及带基点超极限下均封闭（第~\ref{sec:prices-background} 节）。由乘积定理~\cite[定理~\ext{P}{prod:main}]{oaP}，其中有一子类具有同样性质，没有端口，并且在有限 \(\ell_2\) 乘积下也封闭。令 \(\mathcal K\) 表示能在该子类中实现且满足 \(\mathbb E_\alpha D\le1\) 的平方距离表 \(D=d^2\) 的集合。取缩放后放置的乘积可得到各表的凸组合。由连通性，每个表项都有界，逐坐标的极限则可通过带基点超极限实现。因此 \(\mathcal K\) 是紧凸集。

对于有序点对上的任意概率分布 \(\nu\le 2\mu\)，令 \(\beta_\nu=\max_{D\in\mathcal K}\mathbb E_\nu D\)。经缩放可知，对该子类中的每种放置都有 \(\mathbb E_\nu d^2\le\beta_\nu\mathbb E_\alpha d^2\)。用 \(X,Y\) 表示分布 \(\nu\) 下的两个端点。由图熵界~\cite[推论~\ext{P}{cor:graph-entropy}]{oaP} 及图的度数有界，得 \(H_\nu(Y\mid X)\le C(1+\beta_\nu)\)，这里的熵以比特计量。每个有序点对的 \(\nu\) 概率质量至多为 \(2/(rq)\)，所以
\[
    H_\nu(Y\mid X)=H_\nu(X,Y)-H_\nu(X)
               \ge \log_2(rq/2)-\log_2 r=\log_2(q/2).
\]
因此，对所有充分大的 \(q\) 都有 \(\beta_\nu\ge c\log q\)，且关于 \(\nu\) 一致成立。

对 \(\mathcal K\) 与所有分布 \(\nu\le2\mu\) 构成的紧凸集应用双线性极小极大定理，可得同一个 \(D=d^2\in\mathcal K\)，使得对所有这些分布都有 \(\mathbb E_\nu D\ge c\log q\)。若满足 \(D<c\log q\) 的点对的 \(\mu\) 概率质量至少为 \(1/2\)，则将 \(\mu\) 以这些点对为条件所得的分布会违反这一界。因此 \(\mathbb E_\mu d\ge \tfrac12\sqrt{c\log q}\)，而由 Cauchy--Schwarz 不等式有 \(\mathbb E_\alpha d\le1\)。结合价格不等式即得结论。在 \(q\) 有界时，只计算外部比特数即已足够。
\end{proof}

在 \(\operatorname{Tag}(h,h)^{\otimes t}\) 的一个网络中，对每个列表均匀选取一个置换，并将标签固定为该置换。该拓扑实现了 \(r=tNh\) 个有效载荷比特的置换。每个输入比特有 \(q=N\) 个可能的目的地，全都与输入处在同一列表、同一字内比特位置，并且等可能。由引理~\ref{lem:uniform-demand}，\(2V+O(tNh)\ge c\,tNh\sqrt h\)。由于 \(n_w=4N(h+1)\)，对任意固定的 \(0<\sigma<1/2\)，此即证明了式~\eqref{eq:many} 对所有充分大的 \(w\) 成立，且关于 \(t\) 一致。

\subsection{有端口的有限域}

固定一个有端口的有限域。我们使用电路下界论文中的带符号计数张量 \(U_\ell^*\)，它有 \(\ell\) 个外部比特。其表项是一个运行时限为多项式的通用机器的带符号计算路径数在该域中的约化值。多项式规模的验证器给出一份副本的规模为 \(\ell^{O(1)}\) 的精确三价网络。对其猜测串求和，要求每条计算路径中未用的比特均为零，并按结果的符号加权。由精确的批量下界~\cite[定理~\ext{P}{thm:batches}]{oaP}，存在 \(\delta,\xi>0\)，使得对所有充分大的 \(\ell\) 有
\[
       \operatorname{vol}_k\bigl((U_\ell^*)^{\otimes j}\bigr)
              >j\ell^{1+\delta}
       \qquad (1\le j\le 2^{\ell^\xi}).
\]

利用带标签置换的下降过程~\cite[命题~\ext{G}{prop:oracle-descent}]{oaG} 给出一个常数 \(a>0\)。对于所有充分大的 \(h\)，给定一个规模 \(s\) 在 \(N\le s\le N^4\) 范围内、具有 \(e\) 个外部比特的三价网络，该过程用三价门和 \(\operatorname{Tag}(h,h)\) 检测为其张量的 \(d\) 份副本构造一个网络。将每个关联比特（包括检测处的关联比特）都计入规模后，所得规模 \(S\) 满足
\[
            S\le C d\bigl(sh^{-a}+e(\log h)^C\bigr),
           \qquad \log_2 d=O\bigl((\log h)^2\bigr).
\]
这些界关于所给网络一致成立。下降过程所需的规范化在网络上附加相等关系叶节点，并细分自环。其耗费只是一个常数因子，所引的界在每个固定的多项式规模范围内都成立。

选取 \(b>\max\{2,2/\xi\}\) 以及 \(0<\sigma<\min\{a/4,b\delta/4,1/4\}\)。假设存在任意大的 \(w=h+1\)，使某个 \(1\le t\le 2N\) 容许 \(D_w^{\otimes t}\) 的一个体积至多为 \(tn_w w^\sigma\) 的实现。由上述两次列表检测的构造，得到一个规模为 \(O(tNh^{1+\sigma})\)、执行 \(t\) 个带标签置换检测的网络。

将下降网络中的检测每 \(t\) 个分为一批，再用这一实现代替每一批。最后一批中，将未使用的检测固定为恒等置换并取全零有效载荷，此时各检测的值均为一。由于代入的网络实现了完整的乘积张量，无论周围连线怎样收缩，这种代换都保持精确，即使连接的是同一批次内的检测亦然。每个满批次对 \(S\) 贡献 \(3tNh\) 个关联比特。因此，将最后一批的成本单独计入，所得三价网络规模为
\[
                     O\bigl(h^\sigma S+tNh^{1+\sigma}\bigr).
\]

令 \(\ell=\lceil h^b\rceil\)，并从 \(U_\ell^*\) 的多项式规模网络的 \(j=N^2\) 份副本出发。在每次下降时，令 \(x=s/(j\ell)\) 为每个外部比特的规模。新的批次含 \(dj\) 份副本。当 \(h\) 充分大时，其每比特规模满足
\[
 \begin{split}
 x'&\le C'\left(h^{-a+\sigma}x+
        h^\sigma(\log h)^C+\frac{tNh^{1+\sigma}}{dj\ell}\right)\\
   &\le h^{-a/2}x+h^{2\sigma}.
 \end{split}
\]
这里 \(j\ge N^2\)、\(d\ge1\)、\(t\le 2N\)，所以分式一项为 \(O(h^{1+\sigma-b})\)。这就是从 \(N^2\) 份副本出发的原因。

最初，对于某个固定的 \(A\)，有 \(x\le h^A\)。选取固定整数 \(L>2A/a\)。在这 \(L\) 步的每一步中，截至该步的复制因子之积为 \(N^{o(1)}\)。归纳可知，\(j=N^{2+o(1)}\)，且 \(x\) 始终以 \(h\) 的多项式为界，所以当前网络及其规范化网络的规模都处于下降过程的适用范围。迭代最后一个不等式得 \(x=O(h^{2\sigma})\)。此时，因 \(\log_2 j=(2+o(1))h\)，故 \(j\le 2^{\ell^\xi}\)；又由对 \(\sigma\) 的选取，\(h^{2\sigma}=o(\ell^\delta)\)。因此最终的三价网络体积为 \(o(j\ell^{1+\delta})\)，与精确的批量下界矛盾。至此，剩余情形下的定理~\ref{thm:many} 也得到了证明。
